%! Introduction to Functions and Change — Unit 1, Lesson 1.0 — Guided Notes (ANSWER KEY)
% THE in-class document, in the MAIN IDEAS / NOTES shape (LESSON_SHAPE.md §1): the vocabulary
% box, the hook, then ONE two-column guidednotes table. Each row is one idea — a label on the
% left; on the right one or two COMPLETE printed sentences, ONE large pre-drawn display, then a
% two-across grid of numbered problems. Problems run 1..19. No formulas and no notation today.
%   I DO   : the teacher reads the definition, marks up the display, and works the first
%            problem of each grid (1, 5, 9, 13).
%   WE DO  : the class works the rest of each grid, every one cold-called.
%   YOU DO : the last row — problems 16–19, alone, for 14 minutes while the teacher circulates.
% The crux is row 3, MANY-TO-ONE: problem 11 (a repeated output does NOT break the rule) and
% problem 12 (reversing the roles DOES). Problems 17 and 19 transfer it.
% The old read-only "Where This Unit Goes" section is gone: the unit map lives on the unit
% cover, and the Debrief points at it.
%
% PAGE LOCKSTEP with the other half of the pair: the vocabulary write lines -> \ansline, and the
% answer argument of every \pcell and \labelbox filled. Nothing else differs.
% Inside a cell `\\` ENDS THE ROW — break lines with \par. A row cannot break across pages:
% the figure and EACH grid row are separate sub-rows (`\\` then `& ...`, probgrid*).
\documentclass[10pt]{article}
\usepackage{precalculus-article}
\usepackage{precalculus-key}
\newcommand{\TallMath}[1]{$\displaystyle #1\rule[-1.4em]{0pt}{3.2em}$}

\begin{document}

\pageheader{Unit 1, Lesson 1.0}{Guided Notes --- Answer Key}

\vspace{0.12in}

\boxguard
\begin{vocabbox}
Fill in each term as we define it together.
\par\vspace{2pt}
\noindent\textbf{\textcolor{plum}{Input (independent variable):}}\\[1pt]
\ansline{The quantity we choose or track; listed first in an ordered pair.}\\[3pt]
\noindent\textbf{\textcolor{plum}{Output (dependent variable):}}\\[1pt]
\ansline{The quantity that responds to the input; listed second.}\\[3pt]
\noindent\textbf{\textcolor{plum}{Function:}}\\[1pt]
\ansline{A rule that sends each input to exactly one output.}\\[3pt]
\noindent\textbf{\textcolor{plum}{Many-to-one:}}\\[1pt]
\ansline{Two or more inputs share one output. Allowed --- still a function.}\\[3pt]
\noindent\textbf{\textcolor{plum}{One-to-many:}}\\[1pt]
\ansline{One input has two or more outputs. Not allowed --- not a function.}\\[3pt]
\noindent\textbf{\textcolor{plum}{Increasing / decreasing / constant:}}\\[1pt]
\ansline{As the input moves forward, the output rises / falls / holds steady.}
\end{vocabbox}

\vspace{0.1in}

\boxguard
\begin{hookbox}
A vending machine and a claw machine both take your money. Press B4 on the vending
machine and you get the same bag of chips every single time; drop the claw and you
get\ldots{} who knows.

Which machine can you \textit{trust}, and what makes it trustworthy?
\ansline{The vending machine: the same button always gives the same snack. One input, one dependable output --- math calls that a function.}
\end{hookbox}

\small
\begin{guidednotes}
% ── ROW 1: input and output ──────────────────────────────────────────────────
\mainidea[Two quantities]{Input \& Output} &
This unit studies pairs of quantities where one \textbf{responds} to the other. The quantity
we choose or track is the \textbf{input} (the \emph{independent variable}); the quantity that
responds is the \textbf{output} (the \emph{dependent variable}). An ordered pair lists the input
first: (input, output). \textbf{The ``you-control'' test:} ask which quantity you choose --- that
one is the input.

\vspace{4pt}
\begin{center}
\begin{tikzpicture}[scale=0.95, >=stealth]
  \node[draw=plum, very thick, fill=lilac, rounded corners=3pt, minimum width=3.4cm,
        minimum height=1.3cm, align=center, font=\small] (box) at (4.3,0)
        {\textbf{pay-per-minute}\\ \textbf{phone plan}};
  \node[font=\small, align=center] (in) at (0,0) {minutes talked};
  \node[font=\small, align=center] (out) at (8.6,0) {cost of the bill};
  \draw[->, very thick, plum] (in) -- (box);
  \draw[->, very thick, plum] (box) -- (out);
  \node[font=\scriptsize\itshape, above] at (1.55,0.15) {you choose};
  \node[font=\scriptsize\itshape, above] at (7.0,0.15) {it responds};
  \node at (0,-1.05) {\labelbox{2.2cm}{input}};
  \node at (8.6,-1.05) {\labelbox{2.2cm}{output}};
\end{tikzpicture}
\end{center}
\\
& \notesprompt{Name the input and the output. Use the you-control test.}
\begin{probgrid}{|Y|Y|}
\pcell{1}{A lit candle slowly burning down.}{1.3cm}{Input: hours burning. Output: height of the candle.} &
\pcell{2}{A ball dropped from a rooftop.}{1.3cm}{Input: time since the drop. Output: height of the ball.} \\ \hline
\end{probgrid}
\\
& \begin{probgrid}*{|Y|Y|}
\pcell{3}{Pumping gas into a car. Then write one ordered pair that could describe it.}{1.3cm}{Input: gallons pumped. Output: total cost. For example, (10, 35) --- 10 gallons cost \$35.} &
\pcell{4}{A classmate says ``the phone bill decides how many minutes you talk.'' What did they get backward?}{1.3cm}{They swapped the roles. You choose the minutes (input); the bill responds (output).} \\ \hline
\end{probgrid}
\\ \hline
% ── ROW 2: what is a function ────────────────────────────────────────────────
\mainidea[The dependable machine]{Function} &
A rule is a \textbf{function} when each input produces \textbf{exactly one} output --- same
input, same output, every time. An \emph{arrow diagram} shows a rule: inputs on the left, an
arrow from each input to its output on the right. A function has \textbf{exactly one arrow
leaving each input}.

\vspace{4pt}
\begin{center}
\begin{tikzpicture}[>=stealth, font=\small]
  \begin{scope}
    \node[font=\small\bfseries\color{plum}] at (0,1.9) {button};
    \node[font=\small\bfseries\color{plum}] at (2.6,1.9) {snack};
    \foreach \n/\y in {A1/1.1, B4/0.3, C2/-0.5} \node (i\n) at (0,\y) {\n};
    \foreach \n/\y/\k in {chips/1.1/a, pretzels/0.3/b, gum/-0.5/c} \node (o\k) at (2.6,\y) {\n};
    \draw[->, thick, plum] (iA1) -- (oa);
    \draw[->, thick, plum] (iB4) -- (ob);
    \draw[->, thick, plum] (iC2) -- (oc);
    \node at (1.3,-1.45) {\labelbox{2.8cm}{a function}};
  \end{scope}
  \begin{scope}[xshift=6.2cm]
    \node[font=\small\bfseries\color{plum}] at (0,1.9) {claw drop};
    \node[font=\small\bfseries\color{plum}] at (2.6,1.9) {prize};
    \node (d1) at (0,0.7) {drop 1};
    \node (d2) at (0,-0.3) {drop 2};
    \foreach \n/\y/\k in {bear/1.1/a, nothing/0.3/b, ball/-0.5/c} \node (p\k) at (2.6,\y) {\n};
    \draw[->, thick, redacc] (d1) -- (pa);
    \draw[->, thick, redacc] (d1) -- (pb);
    \draw[->, thick, redacc] (d2) -- (pc);
    \node at (1.3,-1.45) {\labelbox{2.8cm}{not a function}};
  \end{scope}
\end{tikzpicture}
\end{center}
\\
& \notesprompt{Function or not? Give the reason in a few words.}
\begin{probgrid}{|Y|Y|}
\pcell{5}{Each student $\to$ their student ID number.}{1.2cm}{Function: each student has exactly one ID number.} &
\pcell{6}{Each month $\to$ the students born that month.}{1.2cm}{Not a function: one month has many students.} \\ \hline
\end{probgrid}
\\
& \begin{probgrid}*{|Y|Y|}
\pcell{7}{Each number $\to$ its double.}{1.2cm}{Function: every number has exactly one double.} &
\pcell{8}{Each shoe size $\to$ a student who wears that size.}{1.2cm}{Not a function: size 9 can point to several students.} \\ \hline
\end{probgrid}
\\ \hline
% ── ROW 3: THE CRUX — many-to-one vs. one-to-many ────────────────────────────
\mainidea[The fine print]{Many-to-One} &
The rule restricts what each \textbf{input} may do. It says nothing about how often an
\textbf{output} may repeat. So \textbf{many-to-one} --- two inputs sharing an output --- is
\textbf{allowed}. \textbf{One-to-many} --- one input with two outputs --- is \textbf{not}.
\textbf{Count the arrows leaving each input, never the arrows arriving at an output.}

\vspace{4pt}
\begin{center}
\begin{tikzpicture}[>=stealth, font=\small]
  \begin{scope}
    \node[font=\small\bfseries\color{plum}] at (0,1.5) {student};
    \node[font=\small\bfseries\color{plum}] at (2.8,1.5) {birthday};
    \node (a1) at (0,0.8) {Ana};
    \node (a2) at (0,0.0) {Ben};
    \node (a3) at (0,-0.8) {Cole};
    \node (b1) at (2.8,0.4) {May 3};
    \node (b2) at (2.8,-0.8) {Jan 9};
    \draw[->, thick, plum] (a1) -- (b1);
    \draw[->, thick, plum] (a2) -- (b1);
    \draw[->, thick, plum] (a3) -- (b2);
    \node[font=\scriptsize\itshape, align=center] at (1.4,-1.5) {Ana and Ben are twins};
    \node at (1.4,-2.3) {\labelbox{3.0cm}{many-to-one: allowed}};
  \end{scope}
  \begin{scope}[xshift=6.2cm]
    \node[font=\small\bfseries\color{plum}] at (0,1.5) {shoe size};
    \node[font=\small\bfseries\color{plum}] at (2.8,1.5) {student};
    \node (s1) at (0,0.4) {9};
    \node (s2) at (0,-0.8) {7};
    \node (t1) at (2.8,0.8) {Dev};
    \node (t2) at (2.8,0.0) {Eli};
    \node (t3) at (2.8,-0.8) {Fay};
    \draw[->, thick, redacc] (s1) -- (t1);
    \draw[->, thick, redacc] (s1) -- (t2);
    \draw[->, thick, redacc] (s2) -- (t3);
    \node[font=\scriptsize\itshape, align=center] at (1.4,-1.5) {two students wear size 9};
    \node at (1.4,-2.3) {\labelbox{3.0cm}{one-to-many: broken}};
  \end{scope}
\end{tikzpicture}
\end{center}
\\
& \notesprompt{Each list is (input, output). Is the rule a function? Say which case it is.}
\begin{probgrid}{|Y|Y|}
\pcell{9}{$(1, 5),\ (2, 5),\ (3, 7)$}{1.2cm}{Function. Inputs 1 and 2 share the output 5 --- many-to-one, allowed.} &
\pcell{10}{$(4, 1),\ (4, 2),\ (6, 3)$}{1.2cm}{Not a function. The input 4 has two outputs, 1 and 2 --- one-to-many.} \\ \hline
\end{probgrid}
\\
& \begin{probgrid}*{|Y|Y|}
\pcell{11}{Two students both scored 80 on a quiz. Maya says, ``So \emph{student $\to$ quiz score} is not a function.'' Is she right?}{1.6cm}{No. Each student still has exactly one score; two inputs sharing the output 80 is many-to-one, which is allowed.} &
\pcell{12}{\emph{Student $\to$ birthday} is a function. Flip it: is \emph{birthday $\to$ student} a function? Use the diagram above.}{1.6cm}{No. The input May 3 would point to two students, Ana and Ben --- one-to-many. Flipping can break a function.} \\ \hline
\end{probgrid}
\\ \hline
% ── ROW 4: a graph tells a story ─────────────────────────────────────────────
\mainidea[A graph tells a story]{Reading Graphs} &
Read a graph \textbf{left to right}, the way the input moves. Where the graph rises, the output
is \textbf{increasing}; where it falls, \textbf{decreasing}; where it is flat, \textbf{constant}.
A flat stretch does not mean nothing is happening --- the input keeps moving; the
\emph{output} holds steady. A cyclist rides away from home, rests, and rides back:
\\
& \begin{center}
\begin{tikzpicture}[x=0.62cm, y=0.46cm]
  \draw[linegray, very thin, step=1] (0,0) grid (10,6);
  \draw[->, thick] (0,0) -- (10.6,0) node[below right] {\small time (min)};
  \draw[->, thick] (0,0) -- (0,6.6) node[above] {\small distance from home (blocks)};
  \foreach \x in {2,4,6,8,10} \node[below] at (\x,0) {\small \x};
  \foreach \y in {1,3,5} \node[left] at (0,\y) {\small \y};
  \draw[very thick, plum] (0,0) -- (4,5) node[midway, above left] {\small\bfseries A};
  \draw[very thick, plum] (4,5) -- (6,5) node[midway, above] {\small\bfseries B};
  \draw[very thick, plum] (6,5) -- (10,0) node[midway, above right] {\small\bfseries C};
\end{tikzpicture}
\end{center}
\vspace{2pt}
A: \labelbox{2.3cm}{increasing}\quad B: \labelbox{2.3cm}{constant}\quad C: \labelbox{2.3cm}{decreasing}
\\
& \notesprompt{Read the story off the graph --- no formula needed.}
\begin{probgrid}{|Y|Y|}
\pcell{13}{How far from home is the rider at $t = 5$? What is the rider doing on segment B?}{1.4cm}{5 blocks. Resting: time keeps moving, but the distance holds steady at 5.} &
\pcell{14}{During which minutes is the distance decreasing? What is the rider doing then?}{1.4cm}{From $t = 6$ to $t = 10$. The rider is riding back toward home.} \\ \hline
\end{probgrid}
\\
& \begin{probgrid}*{|Y|}
\pcell{15}{Is distance from home a function of time? Explain using the graph.}{1.1cm}{Yes. At each time the rider is at exactly one distance --- each input has one output.} \\ \hline
\end{probgrid}
\\ \hline
% ── LAST ROW: YOU DO — alone, while the teacher circulates (14 min) ──────────
\mainidea[You do, alone]{Tables \& Flips} &
A table tells the same story a graph does: watch what the output does as the input moves
forward. Each problem brings its own table; the last two share one runner.
\\
& \notesprompt{On your own. Show your work.}
\begin{probgrid}{|Y|Y|}
\pcell{16}{\tikz[baseline=(t.base)]\node[inner sep=0pt, font=\footnotesize] (t) {\renewcommand{\arraystretch}{1.2}\begin{tabular}{|c|c|c|c|c|c|}\hline Hours & 0 & 1 & 2 & 3 & 4 \\ \hline Height (cm) & 20 & 17 & 14 & 14 & 11 \\ \hline\end{tabular}};\par\vspace{3pt} A candle burns. When is the height decreasing, and when is it constant? Which segment of the bike graph matches the constant stretch?}{2.2cm}{Decreasing from hour 0 to 2 and from hour 3 to 4; constant from hour 2 to 3. It matches segment B.} &
\pcell{17}{\tikz[baseline=(t.base)]\node[inner sep=0pt, font=\footnotesize] (t) {\renewcommand{\arraystretch}{1.2}\begin{tabular}{|c|c|c|c|c|}\hline Hours studied & 0 & 1 & 2 & 3 \\ \hline Quiz score & 60 & 70 & 80 & 80 \\ \hline\end{tabular}};\par\vspace{3pt} Inputs 2 and 3 share the output 80. Is quiz score still a function of hours studied? Explain.}{2.2cm}{Yes. Each number of hours has exactly one score. Two inputs sharing 80 is many-to-one, which is allowed.} \\ \hline
\end{probgrid}
\\
& \begin{probgrid}*{|Y|Y|}
\pcell{18}{A runner jogs out to a bench, rests, and jogs back. She is 0.3 km from the start once on the way out and once on the way back. Is distance a function of time?}{3.2cm}{Yes. At each time she is at exactly one distance. Two times sharing 0.3 km is many-to-one --- allowed.} &
\pcell{19}{Flip it for the same runner: is \emph{time} a function of \emph{distance}? Explain using the 0.3 km moment.}{3.2cm}{No. The input 0.3 km would have two outputs --- two different times. One input with two outputs is not a function.} \\ \hline
\end{probgrid}
\\ \hline
\end{guidednotes}

% ── THE NOTES END AT THE YOU DO ROW; AP Practice and the Homework follow ──────

\end{document}
